\documentclass[11pt,a4paper]{article}

\usepackage[utf8]{inputenc}
\usepackage[T1]{fontenc}
\usepackage[brazil]{babel}
\usepackage{geometry}
\usepackage{booktabs}
\usepackage{xcolor}
\usepackage{helvet}
\usepackage{hyperref}
\usepackage{enumitem}

\renewcommand{\familydefault}{\sfdefault}
\geometry{margin=2.4cm}
\setlength{\parindent}{0pt}
\setlength{\parskip}{0.65em}
\setlist{nosep,leftmargin=1.2cm}

\definecolor{azul}{HTML}{1F4E79}
\definecolor{cinza}{HTML}{F2F4F7}

\hypersetup{
  colorlinks=true,
  linkcolor=azul,
  urlcolor=azul
}

\begin{document}

\begin{titlepage}
\vspace*{2cm}
{\Large Fundação Getulio Vargas\par}
\vspace{0.4cm}
{\Large Metodologia de Survey e Aplicações Eleitorais\par}
\vspace{0.4cm}
{\large Professor Felipe Nunes\par}
\vfill
{\color{azul}\Huge\bfseries Tarefa II\par}
\vspace{0.3cm}
{\color{azul}\LARGE Exercício 2: Elaboração de Questionário\par}
\vfill
{\large Autor: Joao Paulo Zangrandi\par}
{\large 23 de julho de 2026\par}
\end{titlepage}

\section*{Informações gerais}

\begin{center}
\fcolorbox{azul}{cinza}{\parbox{0.9\textwidth}{
\textbf{Tema do questionário:} atitudes políticas e participação eleitoral no Brasil.

\textbf{Modo de coleta:} pesquisa via internet, autoaplicada.

\textbf{Link do Google Forms:} \textbf{a inserir após criação do formulário na conta Google}.
}}
\end{center}

\textbf{Questionário original usado como base:} CESOP-DATAFOLHA/00219, ``Avaliação Sarney / Expectativa Collor II - Cultura e Política III - CEDEC''. A página do CESOP está disponível em: \url{https://www.cesop.unicamp.br/por/banco_de_dados/v/458}. O questionário original está disponível em: \url{https://www.cesop.unicamp.br/vw/1IsfxSa4wNQ_MDA_26dd5_}.

\section*{Mensagem de recrutamento e instruções}

Você está sendo convidado(a) a responder uma pesquisa curta sobre atitudes políticas e participação eleitoral no Brasil. O questionário é anônimo, autoaplicado e leva cerca de 3 minutos. Não existem respostas certas ou erradas; responda de acordo com sua opinião. Sua participação é voluntária e você pode parar a qualquer momento.

\section*{Questionário proposto}

\begin{enumerate}
\item \textbf{Você concorda em participar desta pesquisa anônima?}

Tipo: múltipla escolha. Obrigatória. Pulo: se ``Não concordo'', enviar formulário.
\begin{itemize}
\item Sim, concordo em participar
\item Não concordo em participar
\end{itemize}

\item \textbf{Qual é a sua idade?}

Tipo: múltipla escolha. Obrigatória.
\begin{itemize}
\item 16 ou 17 anos
\item 18 a 29 anos
\item 30 a 44 anos
\item 45 a 59 anos
\item 60 anos ou mais
\item Prefiro não responder
\end{itemize}

\item \textbf{Em qual região do Brasil você mora atualmente?}

Tipo: múltipla escolha. Obrigatória.
\begin{itemize}
\item Norte
\item Nordeste
\item Centro-Oeste
\item Sudeste
\item Sul
\item Moro fora do Brasil
\item Prefiro não responder
\end{itemize}

\item \textbf{Em geral, qual é o seu nível de interesse por política?}

Tipo: escala linear de 1 a 5. Obrigatória. Rotulo 1: nenhum interesse. Rótulo 5: muito interesse.

\item \textbf{Como você avalia o governo federal atual?}

Tipo: múltipla escolha. Obrigatória.
\begin{itemize}
\item Ótimo
\item Bom
\item Regular
\item Ruim
\item Péssimo
\item Não sei avaliar
\end{itemize}

\item \textbf{Se o voto não fosse obrigatório, você votaria na próxima eleição para presidente?}

Tipo: múltipla escolha. Obrigatória.
\begin{itemize}
\item Com certeza votaria
\item Provavelmente votaria
\item Provavelmente não votaria
\item Com certeza não votaria
\item Ainda não sei
\end{itemize}

\item \textbf{Você tem algum partido político de preferência?}

Tipo: múltipla escolha. Obrigatória. Pulo: se ``Sim'', ir para a pergunta 8; caso contrário, ir para a pergunta 9.
\begin{itemize}
\item Sim
\item Não
\item Não sei
\item Prefiro não responder
\end{itemize}

\item \textbf{Qual partido político você prefere?}

Tipo: resposta curta. Obrigatória apenas para quem respondeu ``Sim'' na pergunta 7. Instrução: escreva a sigla ou o nome do partido.

\item \textbf{Na sua opinião, os partidos políticos são importantes para a democracia?}

Tipo: múltipla escolha. Obrigatória.
\begin{itemize}
\item Muito importantes
\item Um pouco importantes
\item Pouco importantes
\item Nada importantes
\item Não sei avaliar
\end{itemize}

\item \textbf{De modo geral, você acha que os políticos brasileiros escutam as demandas da população?}

Tipo: escala linear de 1 a 5. Obrigatória. Rótulo 1: nunca escutam. Rótulo 5: sempre escutam.

\item \textbf{Pensando nas últimas duas semanas, por qual meio você mais se informou sobre política?}

Tipo: múltipla escolha. Obrigatória.
\begin{itemize}
\item Televisão
\item Rádio
\item Jornais ou portais de notícia na internet
\item Redes sociais
\item Conversas com familiares, amigos ou colegas
\item Não me informei sobre política nesse período
\item Outro
\end{itemize}

\item \textbf{Em uma frase, o que mais poderia aumentar sua confiança na política brasileira?}

Tipo: parágrafo. Não obrigatória.
\end{enumerate}

\section*{Justificativa}

O questionário original do CESOP-DATAFOLHA/00219 tinha muitas perguntas presenciais sobre avaliação de governo, voto anterior, preferência partidária, papel dos partidos e percepção sobre políticos. Mantive esse eixo temático, mas adaptei o instrumento para uma aplicação online, curta e autoaplicada. Por isso, reduzi o questionário para 12 perguntas, atualizei referências temporais e substituí perguntas muito dependentes do contexto de 1990 por perguntas mais gerais sobre atitudes políticas atuais.

As escolhas de redação seguem os critérios discutidos por Krosnick (2018) e pela literatura de desenho de questionários: perguntas curtas, linguagem simples, alternativas exclusivas e escalas simétricas. Também evitei perguntas longas, dupla negação e listas extensas demais, pois esse tipo de formulação aumenta a dificuldade cognitiva e pode estimular \textit{satisficing}. Em perguntas de atitude, usei escalas de 5 pontos ou alternativas ordenadas, facilitando o mapeamento entre julgamento e resposta.

Como o modo de coleta é online e autoaplicado, incluí uma mensagem inicial breve, uma pergunta de consentimento e poucos pulos. O pulo principal aparece na preferência partidária: apenas quem declara ter partido de preferência recebe a pergunta aberta sobre qual partido prefere. Essa decisão reduz carga desnecessária para quem não tem preferência partidária e melhora a experiência de resposta, sem copiar literalmente o questionário original.

\end{document}
